\chapter{Siegel's Theorem}
\begin{center}\itshape
A real character might have a zero very close to $1$. Nobody has found one, and nobody can rule it out.\\
Siegel showed that such a zero, if it exists, cannot come too close; but his proof cannot tell how close.\\
It is the most famous ineffective theorem in number theory, and everything about primes in progressions rests on it.
\end{center}

In the \emph{Disquisitiones Arithmeticae} of 1801, Gauss classified binary quadratic forms and tabulated their class numbers. He noticed that the negative discriminants of class number one seem to stop early, and asked, in effect, whether the class number of $\Q(\sqrt{-d})$ tends to infinity with $d$. For imaginary quadratic fields this is
Gauss's class number problem. Dirichlet's class number formula of Volume~13 turns it into analysis: for a prime $p\equiv3\pmod4$, $p>3$,
\[
h(-p)=\frac{\sqrt p}\pi\,L(1,\chi_p),
\]
where $\chi_p$ is the Legendre symbol modulo $p$. So the question is how small $L(1,\chi)$ can be for a real character $\chi$.

The answer came in a strange way. Erich Hecke showed that the generalised Riemann hypothesis would give $L(1,\chi)\gg\frac1{\log q}$. Max Deuring proved in 1933 that if the Riemann hypothesis is \emph{false}, the class number tends to infinity, and Hans Heilbronn extended this in 1934 to the failure of the generalised Riemann
hypothesis. Either way the class number tends to infinity: a proof by cases, one of which is believed never to happen, and which gives no idea of how fast. A year later Carl Ludwig Siegel proved the sharp form, $L(1,\chi)\ge c(\varepsilon)q^{-\varepsilon}$, hence $\log h(-p)\sim\frac12\log p$, with a constant $c(\varepsilon)$ that nobody can compute, to this day.
The class number one problem itself, the list of the nine fields, was settled by Alan Baker and Harold Stark in 1966 and 1967, and Stark then showed that a proof published by Kurt Heegner in 1952, and dismissed at the time, had been essentially correct.

What can go wrong is an \emph{exceptional zero}: a real zero of $L(s,\chi)$ extremely close to $1$, for a real character $\chi$. This chapter proves the zero-free region of Dirichlet $L$-functions, with its one possible exception; shows, after Landau and Page, that exceptional zeros repel each other; proves Siegel's theorem; and draws the consequence
on which all of analytic number theory relies, the Siegel--Walfisz theorem on primes in progressions with a modulus growing like a power of $\log x$.

\begin{reminder}[from Volumes 6 and 13]
A character $\chi$ modulo $q$ is induced by a unique primitive character $\chi^*$ modulo its conductor $q^*\mid q$ (Volume~6, Proposition~3.4), and $L(s,\chi)=L(s,\chi^*)\prod_{p\mid q}(1-\chi^*(p)p^{-s})$. For $\sigma>1$, $L(s,\chi)=\prod_p(1-\chi(p)p^{-s})^{-1}$ and
$-\frac{L'}L(s,\chi)=\sum_n\chi(n)\Lambda(n)n^{-s}$ (Volume~6, Theorem~3.6). For $\chi$ non-principal, $L(s,\chi)$ extends to $\sigma>0$ and $L(1+it,\chi)\ne0$ (Theorems~3.6, 3.8 and~5.8). Jensen's corollary and the local lemma: if $f$ is holomorphic near $|z|\le R$, $f(0)\ne0$ and $|f(z)|\le M|f(0)|$ for $|z|\le R$, $M\ge2$, then $f$ has at most
$\frac{\log M}{\log2}$ zeros $\rho$ in $|z|\le\frac R2$, and $\bigl|\frac{f'}f(z)-\sum_\rho\frac1{z-\rho}\bigr|\le\frac{17\log M}R$ for $|z|\le\frac R4$ (Volume~6, Lemma~6.6). The zeros of $\zeta$ satisfy $\beta<1-\frac c{\log(|t|+2)}$ (Volume~6, Theorem~6.12), and $\psi(x)=x+O(xe^{-c\sqrt{\log x}})$ (Theorem~7.8). Finally, for an odd prime $p\equiv3\pmod4$, $p>3$,
$L(1,\chi_p)=\frac{\pi h(-p)}{\sqrt p}$ (Volume~13, Theorems~11.9 and~11.10, with $K=\Q(\sqrt{-p})$, $r_2=1$, $R_K=1$, $w_K=2$, $|d_K|=p$).
\end{reminder}

\section{Dirichlet \texorpdfstring{$L$}{L}-functions near \texorpdfstring{$s=1$}{s=1}}

Throughout, $q\ge3$, $\chi$ is a character modulo $q$, $s=\sigma+it$, and
\[
\mathcal L=\mathcal L(q,t)=\log\bigl(q(|t|+2)\bigr).
\]
The zeros of $L(s,\chi)$ with $0<\beta<1$ are written $\rho=\beta+i\gamma$. If $\chi$ is induced by $\chi^*$, they are the zeros of $L(s,\chi^*)$ (the extra factors vanish only on $\sigma=0$); in particular the real zeros in $]0,1[$ are the same.

\begin{lemma}[size of $L$]\label{lem:lsize}
Let $\chi$ be non-principal modulo $q$. Then $|\sum_{n\le x}\chi(n)|\le q$ for all $x$, and:
\begin{enumerate}[label=(\alph*)]
\item $|L(s,\chi)|\le\frac{q|s|}\sigma$ for $\sigma>0$;
\item $|L(1,\chi)|\le2+\log q$;
\item if $q\ge8$ and $1-\frac1{\log q}\le\sigma\le1$, then $|L'(\sigma,\chi)|\le7\log^2q$.
\end{enumerate}
\end{lemma}
\begin{proof}
The sum of $\chi$ over a full period is $0$ (Volume~6, Proposition~3.1), so a partial sum is a sum over fewer than $q$ terms of modulus $\le1$. Let $S(u)=\sum_{n\le u}\chi(n)$.
(a) Abel summation gives, for $\sigma>1$ and then by analytic continuation for $\sigma>0$, $L(s,\chi)=s\int_1^\infty S(u)u^{-s-1}\dd u$, of modulus at most $q|s|\int_1^\infty u^{-\sigma-1}\dd u=\frac{q|s|}\sigma$.
(b) $L(1,\chi)=\sum_{n\le q}\frac{\chi(n)}n+\int_q^\infty\frac{S(u)-S(q)}{u^2}\dd u$, and the integral is at most $q\int_q^\infty u^{-2}\dd u=1$ in modulus, the sum at most $1+\log q$.
(c) $-L'(\sigma,\chi)=\sum_n\chi(n)g(n)$ with $g(u)=(\log u)u^{-\sigma}$. Since $q^{1-\sigma}\le q^{1/\log q}=e$, the terms $n\le q$ contribute at most $e\log q\sum_{n\le q}\frac1n\le e\log q(1+\log q)$. For $u\ge q$,
$g'(u)=u^{-\sigma-1}(1-\sigma\log u)\le0$, because $\sigma\log u\ge\log q-1\ge1$. So, by Abel summation, $\bigl|\sum_{n>q}\chi(n)g(n)\bigr|\le2q\,g(q)\le2e\log q$. In total $|L'(\sigma,\chi)|\le e\log q(3+\log q)\le7\log^2q$, as $\log q\ge2$.
\end{proof}

\begin{lemma}[$\frac{L'}L$ near the line $\sigma=1$]\label{lem:lnear}
There is an absolute constant $A_0$ such that for $\chi$ non-principal modulo $q$, $t\in\R$ and $\frac78\le\sigma\le\frac{11}8$,
\[
\Bigl|\frac{L'}L(s,\chi)-\sum_{|\rho-s_0|\le1/2}\frac1{s-\rho}\Bigr|\le A_0\mathcal L,\qquad s_0=\tfrac98+it,
\]
and the sum has at most $A_0\mathcal L$ terms. In particular, for $1<\sigma\le\frac{11}8$, $\operatorname{Re}\bigl(-\frac{L'}L(s,\chi)\bigr)\le A_0\mathcal L$, and $\operatorname{Re}\bigl(-\frac{L'}L(s,\chi)\bigr)\le-\operatorname{Re}\frac1{s-\rho_0}+A_0\mathcal L$ for every zero $\rho_0$ with $|\rho_0-s_0|\le\frac12$.
\end{lemma}
\begin{proof}
Apply the local lemma to $f(z)=L(s_0+z,\chi)$ with $R=1$. For $|z|\le1$, $\operatorname{Re}(s_0+z)\ge\frac18$ and $|s_0+z|\le|t|+3$, so $|f(z)|\le8q(|t|+3)$ by Lemma~\ref{lem:lsize}(a). And $\frac1{L(s_0,\chi)}=\sum_n\mu(n)\chi(n)n^{-s_0}$, so $\frac1{|f(0)|}\le\zeta(\frac98)\le9$, using $\zeta(\sigma)\le1+\int_1^\infty u^{-\sigma}\dd u=\frac\sigma{\sigma-1}$.
So $M=72q(|t|+3)$ will do, and $\log M\le A\mathcal L$. The points $s$ of the statement satisfy $|s-s_0|\le\frac14$. For the last claims: when $\sigma>1>\beta$, every $\operatorname{Re}\frac1{s-\rho}=\frac{\sigma-\beta}{|s-\rho|^2}$ is positive, so dropping terms from $-\sum\operatorname{Re}\frac1{s-\rho}$ can only increase it.
\end{proof}

For the principal character we use $\zeta$. If $\chi_0$ is principal modulo $q$, then $-\frac{L'}L(s,\chi_0)=-\frac{\zeta'}\zeta(s)-\sum_{p\mid q}\frac{\log p\,p^{-s}}{1-p^{-s}}$, and for $\sigma>1$ the last sum has modulus at most $\sum_{p\mid q}\frac{\log p}{p-1}\le2\log q$. And for $1<\sigma\le\frac{11}8$,
\begin{equation}\label{eq:zetanear}
\operatorname{Re}\Bigl(-\frac{\zeta'}\zeta(s)\Bigr)\le\operatorname{Re}\frac1{s-1}+A_0\log(|t|+2):
\end{equation}
for $|t|\ge2$ this follows from Volume~6, Proposition~6.11, by dropping the zeros as above (and $\operatorname{Re}\frac1{s-1}>0$); for $|t|\le2$, $-\frac{\zeta'}\zeta(s)-\frac1{s-1}$ is holomorphic on the compact rectangle $1\le\sigma\le\frac{11}8$, $|t|\le2$, where $\zeta$ has no zero, hence bounded there. Enlarging $A_0$ if necessary, both hold with the same constant.

\section{The zero-free region and the exceptional zero}

\begin{theorem}[Landau, Gronwall, Titchmarsh]\label{thm:lzerofree}
There is an absolute constant $c_1>0$ with the following property. Let $\chi$ be a character modulo $q\ge3$. Then $L(s,\chi)$ has no zero in the region
\[
\sigma\ge1-\frac{c_1}{\mathcal L(q,t)},
\]
except, when $\chi$ is real and non-principal, possibly one zero, which is then real and simple. It is called the \emph{exceptional zero} of $\chi$, and written $\beta_1$.
\end{theorem}

\begin{proof}
For the principal character the zeros are those of $\zeta$, and Volume~6, Theorem~6.12 applies, since $\mathcal L\ge\log(|t|+2)$. Let $\chi$ be non-principal, and $\rho=\beta+i\gamma$ a zero with $\beta\ge\frac58$, so that $|\rho-s_0|\le\frac12$ for $s_0=\frac98+i\gamma$; write $\mathcal L=\mathcal L(q,\gamma)$.

\emph{The 3-4-1 inequality.} For $(n,q)=1$, write $\chi(n)n^{-i\gamma}=e^{i\phi}$: then $3+4\operatorname{Re}(\chi(n)n^{-i\gamma})+\operatorname{Re}(\chi(n)^2n^{-2i\gamma})=3+4\cos\phi+\cos2\phi=2(1+\cos\phi)^2\ge0$. Summing against $\Lambda(n)n^{-\sigma}$, for $\sigma>1$,
\[
\begin{aligned}3\Bigl(-\frac{L'}L(\sigma,\chi_0)\Bigr)&+4\operatorname{Re}\Bigl(-\frac{L'}L(\sigma+i\gamma,\chi)\Bigr)\eqbr+\operatorname{Re}\Bigl(-\frac{L'}L(\sigma+2i\gamma,\chi^2)\Bigr)\ge0.\end{aligned}
\]
Take $1<\sigma\le\frac{11}8$. The first term is at most $3\bigl(\frac1{\sigma-1}+A_0\bigr)$, since its coefficients $\chi_0(n)\Lambda(n)$ are at most those of $-\frac{\zeta'}\zeta$. The second is at most $4\bigl(-\frac1{\sigma-\beta}+A_0\mathcal L\bigr)$ by Lemma~\ref{lem:lnear}.

\emph{Complex characters.} If $\chi^2$ is not principal, the third term is at most $A_0\log(q(2|\gamma|+2))\le2A_0\mathcal L$. Collecting the constants into one $A_1\ge1$,
\[
\frac4{\sigma-\beta}\le\frac3{\sigma-1}+A_1\mathcal L.
\]

\emph{Real characters.} If $\chi^2=\chi_0$ is principal, the third term is at most $\operatorname{Re}\frac1{\sigma-1+2i\gamma}+A_0\log(2|\gamma|+2)+2\log q$ by \eqref{eq:zetanear}, and with the same $A_1$
\[
\frac4{\sigma-\beta}\le\frac3{\sigma-1}+\frac{\sigma-1}{(\sigma-1)^2+4\gamma^2}+A_1\mathcal L.
\]

\emph{The squeeze.} Let $\delta>0$, to be chosen below with $\delta\le\frac1{4A_1}$, and $\sigma=1+\frac\delta{\mathcal L}$. In the complex case, or in the real case with $|\gamma|\ge\frac\delta{\mathcal L}$, the middle term above is at most $\frac{\sigma-1}{4\gamma^2}\le\frac{\mathcal L}{4\delta}$, so in both cases
\[
\frac4{\sigma-\beta}\le\Bigl(\frac3\delta+\frac1{4\delta}+A_1\Bigr)\mathcal L\le\frac{7\mathcal L}{2\delta},\qquad\text{so}\qquad\sigma-\beta\ge\frac{8\delta}{7\mathcal L},\quad1-\beta\ge\frac\delta{7\mathcal L}.
\]

\emph{Real characters, small $\gamma$.} Here the pole of $\zeta$ at $1$ is in the way, and we argue on the real axis. For $1<\sigma'\le\frac{11}8$, the coefficients of $-\frac{\zeta'}\zeta-\frac{L'}L(\cdot,\chi)$ are $\Lambda(n)(1+\chi(n))\ge0$, so $-\frac{L'}L(\sigma',\chi)\ge\frac{\zeta'}\zeta(\sigma')\ge-\frac1{\sigma'-1}-A_0$. Lemma~\ref{lem:lnear} at $t=0$, where $\mathcal L(q,0)\le2\log q$, then gives
\[
\sum_{|\rho'-9/8|\le1/2}\operatorname{Re}\frac1{\sigma'-\rho'}\le\frac1{\sigma'-1}+A_2\log q
\]
with an absolute $A_2$. Let $\lambda=\frac1{2A_2}$, $\kappa=\frac\lambda{10}$, $\sigma'=1+\frac\lambda{\log q}$, and let $B$ be the box $\beta'\ge1-\frac\kappa{\log q}$, $|\gamma'|\le\frac\kappa{\log q}$. A zero $\rho'\in B$ has $\frac\lambda{\log q}\le\sigma'-\beta'\le\frac{1.1\lambda}{\log q}$, so
\[
\operatorname{Re}\frac1{\sigma'-\rho'}=\frac{\sigma'-\beta'}{(\sigma'-\beta')^2+\gamma'^2}\ge\frac{\lambda/\log q}{(1.21+0.01)\lambda^2/\log^2q}\ge\frac{0.81\log q}\lambda .
\]
If $B$ contained two zeros, counted with multiplicity, the left side would be at least $\frac{1.62\log q}\lambda$, while the right side is $(\frac1\lambda+A_2)\log q=\frac{1.5\log q}\lambda$: impossible. So $B$ contains at most one zero. Since $\chi$ is real, $L(\overline s,\chi)=\overline{L(s,\chi)}$, so the zeros are symmetric about the real axis, and $B$ is too: a single zero in $B$ is real, and simple.

\emph{Conclusion.} Choose $\delta=\min(\frac1{4A_1},\kappa)$ and $c_1=\min(\frac\delta8,\frac38)$. Let $\rho$ be a zero with $\beta\ge1-\frac{c_1}{\mathcal L}$ (so $\beta\ge\frac58$). By the squeeze, $\chi$ is real and $|\gamma|<\frac\delta{\mathcal L}$. Then $\mathcal L\ge\log q$ gives $|\gamma|<\frac\kappa{\log q}$ and $\beta\ge1-\frac{c_1}{\log q}\ge1-\frac\kappa{\log q}$: $\rho$ lies in $B$, so it is the one real simple zero of $B$.
\end{proof}

The exceptional zero is the \emph{Siegel zero} of Volume~6. The theorem does not say that it exists; it is believed that it never does, since the generalised Riemann hypothesis puts every zero on $\sigma=\frac12$. Exceptional zeros of different characters repel each other.

\begin{theorem}[Landau, Page]\label{thm:page}
There is an absolute constant $c_2>0$ such that if $\chi_1\ne\chi_2$ are real primitive characters modulo $q_1$ and $q_2$, and $\beta_1$, $\beta_2$ are real zeros of $L(s,\chi_1)$, $L(s,\chi_2)$, then $\min(\beta_1,\beta_2)\le1-\frac{c_2}{\log(q_1q_2)}$.
\end{theorem}
\begin{proof}
The product $\chi_1\chi_2$ is a real character modulo $q_1q_2$, and it is not principal: otherwise $\chi_1(n)=\chi_2(n)$ for all $n$ prime to $q_1q_2$, and $\chi_1$, $\chi_2$ would induce the same character, hence be equal as primitive characters (Volume~6, Proposition~3.4). For $\sigma>1$ the coefficients of
\[
-\frac{\zeta'}\zeta(\sigma)-\frac{L'}L(\sigma,\chi_1)-\frac{L'}L(\sigma,\chi_2)-\frac{L'}L(\sigma,\chi_1\chi_2)
\]
are $\Lambda(n)(1+\chi_1(n))(1+\chi_2(n))\ge0$. Bound the four terms from above for $1<\sigma\le\frac{11}8$ as in the proof of Theorem~\ref{thm:lzerofree}, keeping only $\beta_1$ and $\beta_2$ (we may assume both are at least $\frac58$): with $Q=q_1q_2$ and an absolute $A_3$,
\[
0\le\frac1{\sigma-1}-\frac1{\sigma-\beta_1}-\frac1{\sigma-\beta_2}+A_3\log Q.
\]
If both $\beta_i\ge1-\frac{\kappa'}{\log Q}$, take $\sigma=1+\frac{\lambda'}{\log Q}$ with $\lambda'=\frac1{2A_3}$, $\kappa'=\frac{\lambda'}4$: then $\frac2{\sigma-\beta_i}\cdot\frac1{\log Q}\ge\frac2{1.25\lambda'}=\frac{1.6}{\lambda'}$, while $\frac1{\sigma-1}+A_3\log Q=\frac{1.5}{\lambda'}\log Q$. Contradiction, so $c_2=\min(\kappa',\frac38)$ works.
\end{proof}

\begin{corollary}[Page, 1935]\label{cor:page}
Let $Q\ge3$. Among all real primitive characters of modulus at most $Q$, at most one has an $L$-function with a real zero $\beta>1-\frac{c_2}{2\log Q}$.
\end{corollary}
\begin{proof}
Two such characters $\chi_1\ne\chi_2$ would have $q_1q_2\le Q^2$, and zeros $\beta_i>1-\frac{c_2}{\log(q_1q_2)}$, against Theorem~\ref{thm:page}.
\end{proof}

\section{Siegel's theorem}

How small can $L(1,\chi)$ be? The class number formula gives an easy answer for $\chi_p$: $h(-p)\ge1$, so $L(1,\chi_p)\ge\frac\pi{\sqrt p}$. Siegel replaced $q^{-1/2}$ by $q^{-\varepsilon}$. The proof below is Theodor Estermann's version (1948), as Harold Davenport presents it. Its engine is a single auxiliary function with non-negative
coefficients, built from two real characters.

Let $\chi_1\ne\chi_2$ be real primitive characters modulo $q_1$, $q_2$, and
\[
F(s)=\zeta(s)L(s,\chi_1)L(s,\chi_2)L(s,\chi_1\chi_2),\qquad\lambda=L(1,\chi_1)L(1,\chi_2)L(1,\chi_1\chi_2).
\]
Each $L(1,\chi)$ for a real non-principal $\chi$ is positive: $L(\sigma,\chi)>0$ for $\sigma>1$ by the Euler product, $L$ is continuous and real on $[1,\infty[$, and it does not vanish at $1$. So $\lambda>0$. And $F$ is holomorphic on $\sigma>0$ except for a simple pole at $1$ with residue $\lambda$.

\begin{lemma}\label{lem:Fpos}
$F(s)=\sum_na_nn^{-s}$ for $\sigma>1$, with $a_n\ge0$ and $a_1=1$.
\end{lemma}
\begin{proof}
By the Euler products,
\[
\log F(s)=\sum_p\sum_{k\ge1}\frac{1+\chi_1(p)^k+\chi_2(p)^k+(\chi_1\chi_2)(p)^k}{kp^{ks}},
\]
and the numerator is $(1+\chi_1(p)^k)(1+\chi_2(p)^k)\ge0$. So $\log F$ is a Dirichlet series with non-negative coefficients and no constant term, and so is every power of it; $F=\exp(\log F)=1+\log F+\frac12(\log F)^2+\cdots$ has non-negative coefficients and $a_1=1$.
(All the series converge absolutely for $\sigma>1$, so they may be rearranged.)
\end{proof}

\begin{lemma}\label{lem:Fbound}
There is an absolute constant $C_1\ge1$ such that for $\frac78\le\sigma<1$,
\[
F(\sigma)\ge\frac12-\frac{C_1\lambda}{1-\sigma}(q_1q_2)^{C_1(1-\sigma)}.
\]
\end{lemma}
\begin{proof}
\emph{Taylor coefficients at $2$.} Near $s=2$, $F(s)=\sum_{m\ge0}b_m(2-s)^m$ with $b_m=\frac{(-1)^mF^{(m)}(2)}{m!}=\frac1{m!}\sum_na_n\frac{(\log n)^m}{n^2}\ge0$, and $b_0=F(2)\ge a_1=1$. Also $\frac\lambda{s-1}=\lambda\sum_m(2-s)^m$ for $|s-2|<1$. The function $G(s)=F(s)-\frac\lambda{s-1}$ is holomorphic on $\sigma>0$, so
$G(s)=\sum_m(b_m-\lambda)(2-s)^m$ on the disc $|s-2|<2$.

\emph{Cauchy's estimates.} On the circle $|s-2|=\frac32$: $\sigma\ge\frac12$, $|s|\le\frac72$ and $|s-1|\ge\frac12$. From $\zeta(s)=\frac s{s-1}-s\int_1^\infty\{u\}u^{-s-1}\dd u$ (Volume~6, Exercise~6.6 with $N=1$), $|\zeta(s)|\le7+7=14$. By Lemma~\ref{lem:lsize}(a), $|L(s,\chi_1)|\le7q_1$, $|L(s,\chi_2)|\le7q_2$, $|L(s,\chi_1\chi_2)|\le7q_1q_2$. And $\lambda\le(2+\log q_1q_2)^3$ by Lemma~\ref{lem:lsize}(b), so $|\frac\lambda{s-1}|\le2\lambda$.
Hence $|G|\le M:=5000(q_1q_2)^2+2(2+\log q_1q_2)^3$ on the circle, and $|b_m-\lambda|\le M(\frac23)^m$.

\emph{Summing.} Let $\frac78\le\sigma<1$, so $1<2-\sigma\le\frac98$, and $K\ge1$. Since $\frac23\cdot\frac98=\frac34$, $\sum_{m\ge K}|b_m-\lambda|(2-\sigma)^m\le4M(\frac34)^K$. Using $b_m\ge0$ for $1\le m<K$ and $b_0\ge1$,
\[
F(\sigma)=G(\sigma)+\frac\lambda{\sigma-1}\ge1-\lambda\sum_{m<K}(2-\sigma)^m-\frac\lambda{1-\sigma}-4M\Bigl(\frac34\Bigr)^K=1-\frac{\lambda(2-\sigma)^K}{1-\sigma}-4M\Bigl(\frac34\Bigr)^K,
\]
because $\sum_{m<K}(2-\sigma)^m=\frac{(2-\sigma)^K-1}{1-\sigma}$. Choose $K$ the least integer with $4M(\frac34)^K\le\frac12$. Since $q_1q_2\ge12$, $\log(8M)\le C\log(q_1q_2)$, so $K\le C_1\log(q_1q_2)$ for an absolute $C_1$, and $(2-\sigma)^K\le e^{K(1-\sigma)}\le(q_1q_2)^{C_1(1-\sigma)}$.
\end{proof}

\begin{theorem}[Siegel, 1935]\label{thm:siegel}
For every $\varepsilon>0$ there is $c(\varepsilon)>0$ such that $L(1,\chi)>c(\varepsilon)q^{-\varepsilon}$ for every real primitive character $\chi$ modulo $q\ge3$.
\end{theorem}
\begin{proof}
We may take $\varepsilon\le\frac18$. Let $C_1$ be as in Lemma~\ref{lem:Fbound}, and $I=\bigl[1-\frac\varepsilon{2C_1},1\bigr[\subseteq\bigl[\frac78,1\bigr[$.

\emph{A bad pair $(\chi_1,\beta)$.} We claim that there are a real primitive character $\chi_1$ and a $\beta\in I$ such that $F(\beta)\le0$ for every real primitive $\chi_2\ne\chi_1$.
\begin{itemize}
\item If some real primitive $\chi_1$ has a zero $\beta\in I$, then $L(\beta,\chi_1)=0$, so $F(\beta)=0$ whatever $\chi_2$.
\item Otherwise let $\chi_1$ be any real primitive character, say $\chi_1=\chi_{-4}$ modulo $4$, and $\beta=1-\frac\varepsilon{4C_1}$. Then $\zeta(\beta)<0$: the formula for $\zeta$ above gives $\zeta(\beta)=\frac\beta{\beta-1}-\beta\int_1^\infty\{u\}u^{-\beta-1}\dd u<0$. And $L(\beta,\chi_1)$, $L(\beta,\chi_2)$, $L(\beta,\chi_1\chi_2)$ are positive: each is induced by a real primitive
character $\chi^*$, whose $L$-function is positive at $1$ and has no zero on $[\beta,1]\subseteq I\cup\{1\}$, hence is positive on $[\beta,1]$; and the extra factors $1-\chi^*(p)p^{-\beta}$ are positive. So $F(\beta)<0$.
\end{itemize}

\emph{The bound.} Fix such $\chi_1$ (modulo $q_1$) and $\beta$; they depend on $\varepsilon$ only. For a real primitive $\chi_2\ne\chi_1$ modulo $q_2$, Lemma~\ref{lem:Fbound} and $F(\beta)\le0$ give $\lambda\ge\frac{1-\beta}{2C_1}(q_1q_2)^{-C_1(1-\beta)}\ge\frac{1-\beta}{2C_1}(q_1q_2)^{-\varepsilon/2}$. By Lemma~\ref{lem:lsize}(b),
\[
L(1,\chi_2)=\frac\lambda{L(1,\chi_1)L(1,\chi_1\chi_2)}\ge\frac{(1-\beta)(q_1q_2)^{-\varepsilon/2}}{2C_1(2+\log q_1)(2+\log q_1q_2)}\ge c(\varepsilon)q_2^{-\varepsilon},
\]
since $(2+\log q_1q_2)q_2^{-\varepsilon/2}$ is bounded in terms of $q_1$ and $\varepsilon$. Decreasing $c(\varepsilon)$ to cover $\chi_2=\chi_1$ proves the theorem.
\end{proof}

\begin{think}
Try to compute $c(\varepsilon)$ from the proof. Where exactly do you get stuck?
\end{think}

The constant depends on $\chi_1$ and $\beta$, and which case of the proof occurs is unknown. If no $L$-function has a zero in $I$, everything is explicit. If one does, $c(\varepsilon)$ depends on its modulus $q_1$ and on $1-\beta$, which nobody can bound before finding it. The proof shows that a single bad character would make all the others good, but
it cannot say whether the bad one exists, nor how large its modulus is. Tikao Tatuzawa sharpened this in 1951: for $0<\varepsilon<\frac1{11}$, $L(1,\chi)\ge\frac{0.655\,\varepsilon}{q^\varepsilon}$ for all real primitive $\chi$ modulo $q\ge\max(e^{1/\varepsilon},e^{11.2})$ with at most \emph{one} exception, explicit but unidentified (Exercise~\ref{ex:ineffective}).

\begin{corollary}\label{cor:siegelzero}
For every $\varepsilon>0$ there is $c'(\varepsilon)>0$ such that, for every real non-principal character $\chi$ modulo $q\ge3$, $L(\sigma,\chi)\ne0$ for $\sigma>1-c'(\varepsilon)q^{-\varepsilon}$.
\end{corollary}
\begin{proof}
The real zeros of $L(s,\chi)$ in $]0,1[$ are those of $L(s,\chi^*)$, where $\chi^*$ is primitive modulo $q^*\le q$, and $L(\sigma,\chi^*)>0$ for $\sigma\ge1$; so we may take $\chi$ primitive. Only finitely many characters have $q<8$, and each $L(s,\chi)$ is non-zero on some $[1-\eta_\chi,1]$. For $q\ge8$, let $\beta$ be a real zero. If
$\beta<1-\frac1{\log q}$ we are done as soon as $c'(\varepsilon)q^{-\varepsilon}\le\frac1{\log q}$. Otherwise, by the mean value theorem and Lemma~\ref{lem:lsize}(c), $L(1,\chi)=L(1,\chi)-L(\beta,\chi)\le7(1-\beta)\log^2q$, and Theorem~\ref{thm:siegel} with $\frac\varepsilon2$ gives $1-\beta\ge\frac{c(\varepsilon/2)q^{-\varepsilon/2}}{7\log^2q}\ge c'(\varepsilon)q^{-\varepsilon}$.
\end{proof}

\section{Primes in progressions: the Siegel--Walfisz theorem}

Let
\[
\psi(x,\chi)=\sum_{n\le x}\chi(n)\Lambda(n),\qquad\psi(x;q,a)=\sum_{\substack{n\le x\\ n\equiv a\ (q)}}\Lambda(n).
\]
By orthogonality (Volume~6, Proposition~3.1), for $(a,q)=1$,
\begin{equation}\label{eq:orthpsi}
\psi(x;q,a)=\frac1{\varphi(q)}\sum_{\chi\bmod q}\overline\chi(a)\psi(x,\chi).
\end{equation}
Volume~6 proved $\psi(x;q,a)\sim\frac x{\varphi(q)}$ for each fixed $q$. For sieve methods, and for Bombieri--Vinogradov in Chapter~4, we need it uniformly as $q$ grows with $x$.

\begin{lemma}[truncated Perron formula for characters]\label{lem:perronchi}
For $x\ge3$, $T\ge2$, $a=1+\frac1{\log x}$ and $\chi$ modulo $q$,
\[
\psi(x,\chi)=\frac1{2\pi i}\int_{a-iT}^{a+iT}\Bigl(-\frac{L'}L(s,\chi)\Bigr)\frac{x^s}s\dd s+O\Bigl(\frac{x\log^2x}T+\log x\Bigr),
\]
with an absolute implied constant.
\end{lemma}
\begin{proof}
The proof of Volume~6, Proposition~7.2 integrates $\sum_n\Lambda(n)n^{-s}$ term by term and bounds the difference between each term's integral and its contribution to $\psi_0(x)$ by $\Lambda(n)$ times a quantity depending only on $\frac xn$ and $T$. With coefficients $\chi(n)\Lambda(n)$ of modulus at most $\Lambda(n)$, the same bounds hold
term by term. Finally $\min(1,\frac x{T\langle x\rangle})\le1$, and $\psi(x,\chi)$ and its mid-point version differ by at most $\log x$.
\end{proof}

\begin{theorem}[the prime number theorem for characters]\label{thm:pntchar}
There is an absolute constant $c_3>0$ such that for $x\ge3$ and every non-principal character $\chi$ modulo $q\le e^{\sqrt{\log x}}$,
\[
\psi(x,\chi)=-E_\chi\frac{x^{\beta_1}}{\beta_1}+O\bigl(xe^{-c_3\sqrt{\log x}}\bigr),
\]
where $E_\chi=1$ if $\chi$ has an exceptional zero $\beta_1$ (Theorem~\ref{thm:lzerofree}) and $E_\chi=0$ otherwise. The implied constant is absolute.
\end{theorem}
\begin{proof}
Let $T=e^{\sqrt{\log x}}$ and $\mathcal L_T=\log(q(T+3))$; then $\mathcal L_T\le3\sqrt{\log x}$. Every zero with $|\gamma|\le T+1$, other than $\beta_1$, has $\beta<1-\frac{c_1}{\mathcal L_T}$. Choose $\sigma_1\in\bigl[1-\frac{c_1}{2\mathcal L_T},1-\frac{c_1}{3\mathcal L_T}\bigr]$ at distance at least $\frac{c_1}{12\mathcal L_T}$ from $\beta_1$ (if $\beta_1$ exists),
which is possible since the interval has length $\frac{c_1}{6\mathcal L_T}$.

\emph{The integrand on the new contour.} Let $s=\sigma+it$ with $\sigma_1\le\sigma\le a$, $|t|\le T$. If $\sigma>\frac{11}8$, $|\frac{L'}L(s,\chi)|\le-\frac{\zeta'}\zeta(\frac{11}8)$. Otherwise Lemma~\ref{lem:lnear} applies: its at most $A_0\mathcal L_T$ zeros $\rho\ne\beta_1$ have $|s-\rho|\ge\sigma_1-\beta\ge\frac{c_1}{2\mathcal L_T}$, so they contribute at most $\frac{2A_0}{c_1}\mathcal L_T^2$.
On the left side $\sigma=\sigma_1$, and on the horizontal sides $|t|=T$, the term $\frac1{s-\beta_1}$ is at most $\frac{12\mathcal L_T}{c_1}$. So $|\frac{L'}L(s,\chi)|\le A_4\mathcal L_T^2$ on these three sides.

\emph{Moving the contour.} In the rectangle with vertices $a\pm iT$, $\sigma_1\pm iT$, the function $-\frac{L'}L(s,\chi)\frac{x^s}s$ is holomorphic ($\chi$ is non-principal, $\sigma_1>0$) except at $\beta_1$ if $\beta_1>\sigma_1$, where it has residue $-\frac{x^{\beta_1}}{\beta_1}$. By the residue theorem, the integral of Lemma~\ref{lem:perronchi} equals this residue (if present) plus the integrals over the three other sides, which are at most
\[
A_4\mathcal L_T^2\Bigl(x^{\sigma_1}\int_{-T}^T\frac{\mathrm dt}{|\sigma_1+it|}+2\int_{\sigma_1}^a\frac{x^\sigma}T\dd\sigma\Bigr)\le A_4\mathcal L_T^2\bigl(x^{\sigma_1}\cdot2(2+\log T)+\tfrac{2ex}T\bigr).
\]
If $\beta_1$ exists but $\beta_1\le\sigma_1$, then $\frac{x^{\beta_1}}{\beta_1}\le2x^{\sigma_1}$, and we may add $-\frac{x^{\beta_1}}{\beta_1}$ at this cost.

\emph{Collecting.} $x^{\sigma_1}\le x\exp\bigl(-\frac{c_1\log x}{3\mathcal L_T}\bigr)\le xe^{-\frac{c_1}9\sqrt{\log x}}$, $\frac xT=xe^{-\sqrt{\log x}}$, and $\mathcal L_T^2\log T\le9(\log x)^{3/2}$. With the error $O(\frac{x\log^2x}T+\log x)$ of Lemma~\ref{lem:perronchi}, everything is $O(xe^{-c_3\sqrt{\log x}})$ for any $c_3<\frac{c_1}9$.
\end{proof}

\begin{corollary}[Siegel--Walfisz, 1936]\label{cor:sw}
For every $A>0$ there is $C(A)$ such that for $x\ge3$, $q\le(\log x)^A$ and $(a,q)=1$,
\[
\Bigl|\psi(x;q,a)-\frac x{\varphi(q)}\Bigr|\le C(A)\frac x{(\log x)^A}.
\]
The constant $C(A)$ cannot be computed by the proof.
\end{corollary}
\begin{proof}
We may take $x$ large, so that $(\log x)^A\le e^{\sqrt{\log x}}$. In \eqref{eq:orthpsi}, the principal character gives $\psi(x,\chi_0)=\psi(x)-\sum_{p\mid q}\sum_{p^k\le x}\log p=x+O(xe^{-c\sqrt{\log x}})+O(\log q\log x)$, by Volume~6, Theorem~7.8 (there are at most $\frac{\log q}{\log2}$ primes $p\mid q$, each with at most $\log x$
in total). Every other character gives $O(xe^{-c_3\sqrt{\log x}})$, plus $\frac{x^{\beta_1}}{\beta_1}\le2x^{\beta_1}$ if it has an exceptional zero. By Corollary~\ref{cor:siegelzero} with $\varepsilon=\frac1{2A}$, $1-\beta_1\ge c'q^{-1/(2A)}\ge c'(\log x)^{-1/2}$, so $x^{\beta_1}\le x\exp(-c'\sqrt{\log x})$. Averaging over the $\varphi(q)$ characters,
\[
\psi(x;q,a)=\frac x{\varphi(q)}+O\bigl(xe^{-c''\sqrt{\log x}}\bigr),
\]
with $c''>0$ depending on $A$, and $e^{-c''\sqrt{\log x}}\le C(A)(\log x)^{-A}$. Through $c'(\frac1{2A})$, the constant depends on Siegel's.
\end{proof}

The range $q\le(\log x)^A$ is tiny: the generalised Riemann hypothesis would give $\psi(x;q,a)=\frac x{\varphi(q)}+O(\sqrt x\log^2x)$, meaningful for $q$ up to $\sqrt x/\log^3x$. But Siegel--Walfisz holds for \emph{every} $q$ in its range, and Chapter~4 will show, with the large sieve, that on average over $q\le\sqrt x/(\log x)^B$
the Riemann hypothesis bound holds in a weak form. That is the Bombieri--Vinogradov theorem, and Siegel--Walfisz is one of its two inputs.

\section{Class numbers}

\begin{theorem}[Siegel, 1935]\label{thm:classnumber}
As $p\to\infty$ through primes $p\equiv3\pmod4$,
\[
\log h(-p)\sim\frac12\log p.
\]
More precisely, for every $\varepsilon>0$, $c(\varepsilon)p^{1/2-\varepsilon}\le h(-p)\le\frac1\pi\sqrt p\,(2+\log p)$.
\end{theorem}
\begin{proof}
For $p>3$, $h(-p)=\frac{\sqrt p}\pi L(1,\chi_p)$ by the reminder, where $\chi_p$, the Legendre symbol, is a real primitive character modulo $p$. Lemma~\ref{lem:lsize}(b) gives the upper bound and Theorem~\ref{thm:siegel} the lower bound (with $\frac{c(\varepsilon)}\pi$).
\end{proof}

The same holds for all imaginary quadratic fields, $\log h(d)\sim\frac12\log|d|$ over fundamental discriminants $d<0$, with the same proof: $\zeta_K(s)=\zeta(s)L(s,\chi_d)$ for a real primitive character $\chi_d$ modulo $|d|$, the Kronecker symbol. Volume~13 proves this factorisation for $d=-p$ only; the general case is the same comparison
of Euler factors, using quadratic reciprocity to see that $\chi_d$ is periodic modulo $|d|$. In particular $h(d)\to\infty$, so there are only finitely many imaginary quadratic fields of class number $1$; but the proof gives no bound on the largest one. Heegner, Baker and Stark found them all: $d=-3,-4,-7,-8,-11,-19,-43,-67,-163$. An effective lower
bound, growing like $\log|d|$ up to a factor smaller than any power of $\log|d|$, came only in 1983--1986, from Dorian Goldfeld's reduction to the existence of an elliptic curve with an $L$-function vanishing to order three at its centre, which Benedict Gross and Don Zagier then found.

For a general number field $K$ of degree $n$, the class number formula ties $h_KR_K$ to the residue of $\zeta_K$ at $1$. Richard Brauer proved in 1947, with Siegel's method, the \emph{Brauer--Siegel theorem} (quoted): if $K$ runs through number fields that are Galois over $\Q$, with $\frac n{\log|d_K|}\to0$, then
\[
\log(h_KR_K)\sim\log\sqrt{|d_K|}.
\]
For real quadratic fields this says that $h_KR_K$, not $h_K$, grows like $\sqrt{d_K}$; Gauss's conjecture that infinitely many real quadratic fields have class number $1$ is still open.

\section{Exercises}

\begin{exercise}[$\star$]\label{ex:hge1}
Using only $h(-p)\ge1$, show that for primes $p\equiv3\pmod4$, $p>3$, $L(1,\chi_p)\ge\frac\pi{\sqrt p}$, and deduce that every real zero $\beta$ of $L(s,\chi_p)$ satisfies $1-\beta\ge\frac\pi{7\sqrt p\log^2p}$ for $p\ge11$. Why is this bound effective, and why is it much weaker than Corollary~\ref{cor:siegelzero}?
\end{exercise}

\begin{exercise}[$\star$]\label{ex:ineffective}
Suppose a computer search found a real primitive $\chi_1$ modulo $q_1$ with a zero $\beta\in\bigl[1-\frac\varepsilon{2C_1},1\bigr[$. Write down an explicit value of $c(\varepsilon)$ that the proof of Theorem~\ref{thm:siegel} would then give. Suppose instead that one proved that no such character exists. What would the proof give? Explain why neither outcome is available today.
\end{exercise}

\begin{exercise}[$\star\star$]\label{ex:leastprime}
Use Corollary~\ref{cor:sw} to show that for every $\varepsilon>0$ there is $C(\varepsilon)$ such that every progression $a\bmod q$ with $(a,q)=1$ contains a prime $p\le C(\varepsilon)\exp(q^\varepsilon)$. (Linnik's theorem, $p\le q^L$ for an absolute $L$, is much deeper.)
\end{exercise}

\begin{exercise}[$\star\star$]\label{ex:onebad}
Show that for every $q\ge3$, at most one real non-principal character modulo $q$ has an exceptional zero in the sense of Theorem~\ref{thm:lzerofree}, provided $c_1$ is replaced by $\min(c_1,\frac{c_2}2)$.
\end{exercise}

\begin{exercise}[$\star\star$]\label{ex:hupper}
Let $p\equiv3\pmod4$ be a prime, $p>3$, and let $T(p)$ be the number of quadratic residues minus the number of non-residues in $\left]0,\frac p2\right[$. Volume~13, Exercise~11.1 shows $h(-p)=\frac{T(p)}{2-\chi_p(2)}$.
\begin{enumerate}
\item Deduce that $\left]0,\frac p2\right[$ always contains more residues than non-residues.
\item Deduce from Theorem~\ref{thm:classnumber} that $T(p)=p^{\frac12+o(1)}$.
\item Rederive the upper bound $h(-p)=O(\sqrt p\log p)$ from Pólya--Vinogradov (Volume~6, Theorem~11.11), without $L(1,\chi_p)$.
\end{enumerate}
\end{exercise}

\begin{exercise}[$\star\star\star$]\label{ex:bias}
Suppose a real primitive character $\chi_1$ modulo $q$ had an exceptional zero $\beta_1=1-\frac1{\log^2q}$. Show that for $x=\exp(C\log^2q)$ with an absolute $C$ large enough, and $q$ large,
\[
\psi(x;q,a)=\frac x{\varphi(q)}\Bigl(1-\chi_1(a)e^{-C}+O\bigl(\tfrac1{\log q}\bigr)\Bigr)+O\Bigl(\frac{x}{\varphi(q)}q^{-1}\Bigr)
\]
uniformly in $(a,q)=1$: primes would avoid the residues of $\chi_1$, at a scale where the prime number theorem for progressions should already hold.
\end{exercise}

\section{Hints}

\begin{hint}{ex:hge1}
Use the class number formula, then the proof of Corollary~\ref{cor:siegelzero}.
\end{hint}

\begin{hint}{ex:ineffective}
In the first case, follow the constants through Lemma~\ref{lem:Fbound}; in the second, the proof uses $\chi_1=\chi_{-4}$ and an explicit $\beta$.
\end{hint}

\begin{hint}{ex:leastprime}
Take $x=\exp(q^{1/A})$ with $A=\frac1\varepsilon$, and compare $\psi(x;q,a)$ with the contribution of prime powers.
\end{hint}

\begin{hint}{ex:onebad}
A real character modulo $q$ with an exceptional zero is induced by a real primitive character of modulus at most $q$. Apply Corollary~\ref{cor:page} with $Q=q$.
\end{hint}

\begin{hint}{ex:hupper}
$h(-p)\ge1$ and $1\le2-\chi_p(2)\le3$. For (c), $T(p)$ is a character sum over an interval.
\end{hint}

\begin{hint}{ex:bias}
Use Theorem~\ref{thm:pntchar} for every character, Corollary~\ref{cor:page} to rule out a second exceptional zero, and choose $C$ so that $e^{-c_3\sqrt C\log q}\le q^{-2}$.
\end{hint}

\section{Solutions}

\begin{solution}{ex:hge1}
\idea The class number is a positive integer: that is a free, effective lower bound for $L(1,\chi_p)$, and a lower bound on $L(1,\chi)$ pushes real zeros away from $1$.
\writeup By the class number formula $L(1,\chi_p)=\frac{\pi h(-p)}{\sqrt p}\ge\frac\pi{\sqrt p}$. Let $\beta$ be a real zero. If $\beta<1-\frac1{\log p}$, then $1-\beta>\frac1{\log p}\ge\frac\pi{7\sqrt p\log^2p}$. Otherwise, as in the proof of Corollary~\ref{cor:siegelzero}, $\frac\pi{\sqrt p}\le L(1,\chi_p)\le7(1-\beta)\log^2p$ (Lemma~\ref{lem:lsize}(c), valid for $p\ge11$).
Every constant here is explicit. But $p^{-1/2}$ is a far weaker saving than $p^{-\varepsilon}$: it does not even exclude zeros at distance $p^{-1/2+o(1)}$ from $1$, which would wreck the prime number theorem for progressions modulo $p$ below $x=e^{\sqrt p}$. Siegel's theorem buys the exponent $\varepsilon$ at the price of effectivity.
\end{solution}

\begin{solution}{ex:ineffective}
\idea Every quantity in Lemma~\ref{lem:Fbound} is explicit; the only unknowns are which case occurs and, in the first case, $q_1$ and $\beta$.
\writeup In the first case the proof gives, for $\chi_2\ne\chi_1$,
\[
L(1,\chi_2)\ge\frac{(1-\beta)q_1^{-\varepsilon/2}}{2C_1(2+\log q_1)(2+\log q_1q_2)}q_2^{-\varepsilon/2},
\]
and $(2+\log q_1q_2)\le(2+\log q_1)\bigl(1+\frac2\varepsilon q_2^{\varepsilon/2}\bigr)$ (as $\log q_2\le\frac2\varepsilon q_2^{\varepsilon/2}$), so $c(\varepsilon)=\frac{\varepsilon(1-\beta)q_1^{-\varepsilon/2}}{6C_1(2+\log q_1)^2}$ works (and $L(1,\chi_1)$ itself is a known number). Once the bad character is found, everything is explicit. In the second case $\chi_1=\chi_{-4}$, $q_1=4$, $\beta=1-\frac\varepsilon{4C_1}$, and
the same formula gives an explicit constant. But to know that we are in the second case, we would have to prove that no real $L$-function has a zero in $I$, a statement close to the generalised Riemann hypothesis for real zeros. And to be in the first case usefully we would need to find the bad character, which, if it exists, is at an unknown and presumably astronomical modulus. Since the
proof cannot tell which case holds, it cannot produce a number. This is why Siegel--Walfisz, and every result that uses it (Bombieri--Vinogradov included), has an ineffective constant.
\end{solution}

\begin{solution}{ex:leastprime}
\idea Siegel--Walfisz finds $\gg\frac x{\varphi(q)}$ prime powers in the progression as soon as $q$ is a fixed power of $\log x$, that is once $x\ge\exp(q^{\varepsilon})$. Prime powers $p^k$ with $k\ge2$ are far too few to account for them.
\writeup Let $A=\frac2\varepsilon$ and $x=\exp(q^{\varepsilon})$, so that $(\log x)^A=q^2\ge q$ and Corollary~\ref{cor:sw} applies. Its error is $C(A)\frac x{(\log x)^A}=C(A)\frac x{q^2}\le\frac x{2q}$ for $q\ge2C(A)$, and $\frac x{\varphi(q)}\ge\frac xq$, so $\psi(x;q,a)\ge\frac x{2q}$.
The prime powers $p^k\le x$ with $k\ge2$ contribute at most $\psi(x)-\vartheta(x)=O(\sqrt x\log x)$, which is less than $\frac x{4q}$ for $q$ large, since $q=(\log x)^{1/\varepsilon}$ is tiny compared with $\frac{\sqrt x}{\log x}$. So some prime $p\le\exp(q^\varepsilon)$ is $\equiv a\pmod q$, for $q\ge q_0(\varepsilon)$. The finitely many moduli $q<q_0(\varepsilon)$, each with finitely many classes, are covered by Dirichlet's theorem after enlarging $C(\varepsilon)$. Like $C(A)$, the constant $C(\varepsilon)$ cannot be computed by this proof.
\end{solution}

\begin{solution}{ex:onebad}
\idea Two exceptional characters modulo the same $q$ would be two real primitive characters of modulus at most $q$ with zeros very close to $1$, which Page forbids.
\writeup Let $\chi\ne\chi'$ be real non-principal characters modulo $q$ with zeros $\beta,\beta'\ge1-\frac{\min(c_1,c_2/2)}{\mathcal L(q,0)}$. Since $\mathcal L(q,0)=\log2q\ge\log q$, we get $\beta,\beta'\ge1-\frac{c_2}{2\log q}$. They are zeros of $L(s,\chi^*)$ and $L(s,\chi'^*)$, where the induced primitive characters $\chi^*\ne\chi'^*$ (different characters modulo $q$ are induced by
different primitive characters, by Volume~6, Proposition~3.4) are real with moduli at most $q$. Corollary~\ref{cor:page} with $Q=q$ forbids this.
\end{solution}

\begin{solution}{ex:hupper}
\idea Volume~13 turned the class number into a count of residues on half the circle. Positivity of $h$ then becomes a statement about residues, and Siegel's theorem becomes a statement about how lopsided the count is.
\writeup (a) Since $h(-p)\ge1$ and $2-\chi_p(2)\in\{1,3\}$, $T(p)=(2-\chi_p(2))h(-p)\ge1>0$. For $p=7$: the residues $1,2$ against the non-residue $3$. For $p=23$: $h(-23)=3$ and $\chi_{23}(2)=1$, so $T=3$; the residues $1,2,3,4,6,8,9$ beat the non-residues $5,7,10,11$ by $7-4=3$.
(b) $h(-p)\le T(p)\le3h(-p)$, so $\log T(p)=\log h(-p)+O(1)=(\frac12+o(1))\log p$ by Theorem~\ref{thm:classnumber}.
(c) $T(p)=\sum_{0<a<p/2}\chi_p(a)$ is a sum of a non-principal character over an interval, so Pólya--Vinogradov gives $|T(p)|\le C\sqrt p\log p$, and $h(-p)\le T(p)$. The lower bound in (b) is the hard half: no bound on character sums can show they are \emph{large}, and the positivity in (a) comes from $h\ge1$, that is, from $L(1,\chi_p)>0$. This route to (a), through $L(1,\chi_p)$, goes back to Dirichlet's class number formula of 1839--1840.
\end{solution}

\begin{solution}{ex:bias}
\idea Every character except $\chi_1$ contributes only a power saving at this $x$, while $\chi_1$ contributes $x^{\beta_1}=xe^{-C}$, a fixed fraction of $x$, with the sign $\chi_1(a)$.
\writeup Let $x=\exp(C\log^2q)$, so $\sqrt{\log x}=\sqrt C\log q$ and $q\le e^{\sqrt{\log x}}$. By Theorem~\ref{thm:pntchar}, each non-principal $\chi\ne\chi_1$ has $\psi(x,\chi)=-E_\chi\frac{x^{\beta_\chi}}{\beta_\chi}+O(xq^{-c_3\sqrt C})$. By Corollary~\ref{cor:page} (with $Q=q$, and the induced primitive characters as in Exercise~\ref{ex:onebad}), no real character $\chi\ne\chi_1$ modulo $q$ has a real zero above
$1-\frac{c_2}{2\log q}$ for $q$ large (as $\beta_1=1-\frac1{\log^2q}$ is above it); so if $E_\chi=1$, $x^{\beta_\chi}\le x\exp(-\frac{c_2C\log^2q}{2\log q})=xq^{-c_2C/2}$. For $\chi_1$, $\frac{x^{\beta_1}}{\beta_1}=x e^{-C}(1+O(\frac1{\log^2q}))$. The principal character gives $x+O(xq^{-c\sqrt C})$ by Volume~6, Theorem~7.8 (and $O(\log q\log x)$). Choose $C$ so large that
$c_3\sqrt C\ge2$, $c\sqrt C\ge2$ and $\frac{c_2C}2\ge2$. Then \eqref{eq:orthpsi} gives
\[
\psi(x;q,a)=\frac1{\varphi(q)}\Bigl(x-\chi_1(a)xe^{-C}\bigl(1+O(\tfrac1{\log^2q})\bigr)+\varphi(q)\,O(xq^{-2})\Bigr),
\]
which is the claim (with room to spare in the first error term). Since $\chi_1(a)=1$ for half the residues and $-1$ for the other half, the residues would get $(1-e^{-C})$ and $(1+e^{-C})$ times their fair share of primes up to $x=q^{C\log q}$. Siegel's theorem says that such a zero cannot exist for $q$ large in terms of $C$ (it would need $1-\beta_1\ge c'(\varepsilon)q^{-\varepsilon}$), but it cannot say from which $q$ on.
\end{solution}

\section*{Sources and further reading}

\begin{itemize}
\item \emph{Siegel's theorem.} C.~L.~Siegel, ``\"Uber die Classenzahl quadratischer Zahlk\"orper'', \emph{Acta Arith.}\ 1 (1935), 83--86. T.~Estermann, ``On Dirichlet's $L$ functions'', \emph{J.\ London Math.\ Soc.}\ 23 (1948), 275--279. D.~Goldfeld, ``A simple proof of Siegel's theorem'', \emph{Proc.\ Nat.\ Acad.\ Sci.\ USA} 71 (1974), 1055. T.~Tatuzawa, ``On a theorem of Siegel'', \emph{Japan.\ J.\ Math.}\ 21 (1951), 163--178.
\item \emph{Exceptional zeros and progressions.} A.~Page, ``On the number of primes in an arithmetic progression'', \emph{Proc.\ London Math.\ Soc.}\ 39 (1935), 116--141. A.~Walfisz, ``Zur additiven Zahlentheorie II'', \emph{Math.\ Z.}\ 40 (1936), 592--607. H.~Davenport, \emph{Multiplicative Number Theory}, 3rd ed., revised by H.~L.~Montgomery, Springer GTM 74, 2000, Chapters 14, 20--22, which this chapter follows.
\item \emph{The class number problem.} D.~Goldfeld, ``Gauss's class number problem for imaginary quadratic fields'', \emph{Bull.\ Amer.\ Math.\ Soc.}\ 13 (1985), 23--37: the history told here, from Gauss to Deuring, Heilbronn, Siegel, Heegner, Baker, Stark, Goldfeld, Gross and Zagier. H.~Heilbronn, ``On the class-number in imaginary quadratic fields'', \emph{Quart.\ J.\ Math.\ Oxford} 5 (1934), 150--160.
K.~Heegner, ``Diophantische Analysis und Modulfunktionen'', \emph{Math.\ Z.}\ 56 (1952), 227--253. A.~Baker, ``Linear forms in the logarithms of algebraic numbers'', \emph{Mathematika} 13 (1966), 204--216. H.~M.~Stark, ``A complete determination of the complex quadratic fields of class-number one'', \emph{Michigan Math.\ J.}\ 14 (1967), 1--27. B.~H.~Gross and D.~B.~Zagier, ``Heegner points and derivatives of $L$-series'',
\emph{Invent.\ Math.}\ 84 (1986), 225--320.
\item \emph{Brauer--Siegel.} R.~Brauer, ``On the zeta-functions of algebraic number fields'', \emph{Amer.\ J.\ Math.}\ 69 (1947), 243--250.
\end{itemize}
